% Generated draft: requires mathematical and source review.
\documentclass[11pt]{article}
\usepackage[margin=25mm]{geometry}
\usepackage{fontspec}
\setmainfont{Latin Modern Roman}
\usepackage{amsmath,amssymb,mathtools,unicode-math}
\setmathfont{Latin Modern Math}
\usepackage{enumitem,needspace,xurl,hyperref}
\setlength{\parindent}{0pt}
\setlength{\parskip}{5pt}
\setlength{\emergencystretch}{3em}
\newcommand{\N}{\mathbb{N}}
\newcommand{\Z}{\mathbb{Z}}
\newcommand{\Q}{\mathbb{Q}}
\newcommand{\R}{\mathbb{R}}
\newcommand{\C}{\mathbb{C}}
\newcommand{\F}{\mathbb{F}}
\newcommand{\A}{\mathbb{A}}
\newcommand{\PP}{\mathbb{P}}
\newcommand{\E}{\mathbb{E}}
\newcommand{\eps}{\varepsilon}
\newcommand{\dd}{\,\mathrm{d}}
\begin{document}
\section*{UnsolvedMath ID: 30001687}
\textbf{Source problem code:} catalogue entry
\subsection{Definitions and mathematical statement}
The supplied problem concerns the thickness of spectra of Fibonacci Hamiltonians. It does not specify the operator, its Hilbert space, potential normalization, phase convention, or coupling range. Consequently, the following is a conditional formulation of the intended question, requiring those choices to be supplied before it defines a particular spectral problem.

For an intended family of bounded self-adjoint operators \(H_\lambda\), write \(\Sigma_\lambda=\sigma(H_\lambda)\), where \(\sigma(H)=\{E\in\R:H-E\operatorname{Id}\text{ has no bounded inverse}\}\). The standard interpretation of thickness is Newhouse thickness, defined as follows. Let \(K\subset\R\) be a nonempty compact perfect nowhere dense set, and let \(J\) be its convex hull. A presentation \(\mathcal P=(U_n)_{n\geq1}\) enumerates the bounded connected components of \(J\setminus K\). For an endpoint \(u\) of \(U_n\), let \(B_{\mathcal P}(u)\) be the connected component of \(J\setminus\bigcup_{j=1}^{n}U_j\) containing \(u\). With \(|\cdot|\) denoting interval length, define
\[
\tau(K)=\sup_{\mathcal P}\inf_{n\geq1}\inf_{u\in\partial U_n}\frac{|B_{\mathcal P}(u)|}{|U_n|}.
\]
This interpretation requires confirmation that the intended spectra are Cantor sets and that this is the intended notion of thickness.

For a specified positive coupling range \(I\subset(0,\infty)\) accumulating at zero, the monotonicity question can be expressed as
\[
\forall\lambda_1,\lambda_2\in I,\qquad
\lambda_1<\lambda_2\ \Longrightarrow\ 
\tau(\Sigma_{\lambda_1})\geq\tau(\Sigma_{\lambda_2}).
\]
Here positive coupling and nonincreasing monotonicity are explicit conditional interpretations: the source does not state the coupling range or whether strict decrease is intended.

The weak-coupling phrase ``behaves like \(1/\lambda\)'' also requires qualification. Comparability means
\[
\exists c,C,\delta>0\ \forall\lambda\in I\cap(0,\delta),\qquad
\frac{c}{\lambda}\leq\tau(\Sigma_\lambda)\leq\frac{C}{\lambda}.
\]
Asymptotic equivalence would instead mean
\[
\lim_{\substack{\lambda\downarrow0\\\lambda\in I}}\lambda\tau(\Sigma_\lambda)=1,
\]
which is a stronger assertion. The supplied dated assessment reports comparability as established by weak-coupling work, while reporting that monotone decrease was conjectured and was not established by the located papers. These are supplied curation claims, not independently verified conclusions. The broader request for careful computations of spectral structure specifies no additional mathematical target.
\subsection{Short English statement}
For the intended Fibonacci Hamiltonian family, determine whether spectral thickness decreases as the coupling increases, and clarify its growth near zero coupling. The supplied assessment reports bounds comparable to \(1/\lambda\), while leaving monotonicity unestablished by the located papers. The operator conventions, coupling range, strictness of monotonicity, and meaning of ``behaves like'' require clarification.
\subsection{Sources}
\begin{itemize}
\item Local UnsolvedMath registry, \texttt{data/unsolvedmath/problems.json}: permanent ID \(30001687\), title ``Spectral Thickness of Fibonacci Hamiltonians,'' exact source problem code \texttt{catalogue entry}; collection: Oberwolfach Reports Open Problems. The record supplies the proposed year 2011 but identifies no individual proposer.
\item Supplied report citation: \emph{Mini-Workshop: Dynamics of Trace Maps and Applications to Spectral Theory} (2011), Oberwolfach Reports, pp.~141--167, DOI \texttt{10.4171/owr/2011/03}. Supplied source URL and literature-source label ``original OWR report'': \url{https://doi.org/10.4171/owr/2011/03}.
\item Supplied literature-source label ``Damanik--Gorodetski weak coupling results'': \url{https://arxiv.org/abs/1001.2552}. The record supplies no fuller bibliographic citation for this entry.
\item Literature assessment checked at \texttt{2026-08-21}: supplied dated curation, not independently verified. Its status label is \texttt{partially\_solved}; it reports thickness comparable to \(\lambda^{-1}\) at weak coupling and says that monotone decrease was conjectured and is not established by the located papers. The background explicitly requests independent verification of time-sensitive status and cited claims.
\item The supplied statement-status label is \texttt{corrected\_verified}. Its accompanying verification note records only the addition of mathematical delimiters to TeX notation without changing mathematical content; it does not supply mathematical verification.
\end{itemize}
\end{document}
